% ===== PRÉAMBULE CANONIQUE — à coller VERBATIM en tête de CHAQUE .tex =====
\documentclass[11pt,a4paper]{article}
\usepackage[utf8]{inputenc}\usepackage[T1]{fontenc}\usepackage{lmodern}
\usepackage[french]{babel}
\usepackage{amsmath,amssymb,mathtools}\usepackage{icomma}
\usepackage{siunitx}\sisetup{locale=FR}  % \qty{}{}, \si{}, \ang{}
\usepackage{xcolor}\usepackage{colortbl}
\usepackage{tikz}\usepackage{pgfplots}\pgfplotsset{compat=1.18}
\usepackage[siunitx,europeanresistors]{circuitikz} % circuits (élec.)
\usepackage[most]{tcolorbox}\usepackage{enumitem}\usepackage{array,booktabs}
\usepackage[a4paper,margin=2.2cm]{geometry}
\usetikzlibrary{arrows.meta,calc,angles,quotes,positioning,patterns,%
  backgrounds,shapes.geometric,decorations.pathmorphing,decorations.markings,intersections}
\definecolor{primary}{RGB}{222,49,99}\definecolor{primarydark}{RGB}{150,22,55}
\definecolor{defcol}{RGB}{0,114,178}\definecolor{methcol}{RGB}{52,92,148}
\definecolor{excol}{RGB}{0,135,90}\definecolor{attncol}{RGB}{200,40,55}
\tcbset{boxbase/.style={enhanced,breakable,boxrule=0.9pt,arc=2.5pt,
  left=3mm,right=3mm,top=2.3mm,bottom=2.3mm,fonttitle=\bfseries,coltitle=white}}
% --- boîtes TUTORIEL (noms EXACTS, ne pas renommer) ---
\newtcolorbox{cle}[1][]{boxbase,colback=defcol!6,colframe=defcol,
  title={\textsc{À retenir}\ifx\relax#1\relax\relax\else\textmd{ : \itshape #1}\fi}}
\newtcolorbox{methode}[1][]{boxbase,colback=methcol!7,colframe=methcol,sharp corners,
  title={\textsc{Méthode}\ifx\relax#1\relax\relax\else\textmd{ --- \itshape #1}\fi}}
\newtcolorbox{exemplebox}[1][]{boxbase,colback=excol!5,colframe=excol,
  title={\textsc{Exemple}\ifx\relax#1\relax\relax\else\textmd{ : \itshape #1}\fi}}
\newtcolorbox{piege}[1][]{boxbase,colback=attncol!6,colframe=attncol,
  title={\textsc{Piège}\ifx\relax#1\relax\relax\else\textmd{ : \itshape #1}\fi}}
\newtcolorbox{remarque}[1][]{boxbase,colback=black!4,colframe=black!45,
  title={\textsc{Remarque}\ifx\relax#1\relax\relax\else\textmd{ : \itshape #1}\fi}}
% --- boîtes EXERCICE / CORRIGÉ (noms EXACTS, auto-comptées) ---
\newtcolorbox[auto counter]{exo}[1][]{boxbase,colback=primary!4,colframe=primary,
  title={\textsc{Exercice}~\thetcbcounter\ifx\relax#1\relax\relax\else\textmd{ --- \itshape #1}\fi}}
\newtcolorbox[auto counter]{corrige}[1][]{boxbase,colback=excol!4,colframe=excol,
  title={\textsc{Corrigé}~\thetcbcounter\ifx\relax#1\relax\relax\else\textmd{ --- \itshape #1}\fi}}
\newcommand{\vv}[1]{\vec{#1}}\newcommand{\ds}{\displaystyle}
\newcommand{\R}{\mathbb{R}}\newcommand{\N}{\mathbb{N}}\newcommand{\Z}{\mathbb{Z}}
% =========================================================================
\begin{document}
\usetikzlibrary{babel}
\begin{exo}[filament : résistance à chaud et à froid]
Une lampe à incandescence porte les indications $\qty{200}{\volt}$ et $\qty{100}{\watt}$ (valeurs
lampe \emph{allumée}). La résistance du filament allumé est \textbf{dix fois} plus grande que celle du
filament éteint.
\begin{enumerate}[label=\alph*)]
  \item Calcule la résistance du filament lorsque la lampe est allumée.
  \item En déduire sa résistance lorsqu'elle est éteinte.
  \item Pourquoi la résistance dépend-elle de l'état (allumé/éteint) ?
\end{enumerate}
\end{exo}

\begin{exo}[série ou parallèle : quel montage chauffe le plus ?]
Deux résistances \emph{identiques} $R$ sont montées, sous la même tension $U$ et pendant la même
durée $t$, soit en série (circuit 1), soit en parallèle (circuit 2).
\begin{center}
\begin{minipage}{0.46\textwidth}\centering
\begin{circuitikz}[scale=0.78,transform shape]
  \draw (0,0) to[battery1,l_=$U$] (0,2) to[R,l=$R$] (2.2,2) to[R,l=$R$] (4.4,2) -- (4.4,0) -- (0,0);
\end{circuitikz}\\[1pt]{\footnotesize Circuit 1 — série}
\end{minipage}\hfill
\begin{minipage}{0.46\textwidth}\centering
\begin{circuitikz}[scale=0.78,transform shape]
  \draw (0,0) to[battery1,l_=$U$] (0,2) -- (3.2,2);
  \draw (0,0) -- (3.2,0);
  \draw (1.7,2) to[R,l=$R$] (1.7,0);
  \draw (3.2,2) to[R,l=$R$] (3.2,0);
\end{circuitikz}\\[1pt]{\footnotesize Circuit 2 — parallèle}
\end{minipage}
\end{center}
\begin{enumerate}[label=\alph*)]
  \item Exprime la chaleur dégagée $W_1$ (série) puis $W_2$ (parallèle) en fonction de $U$, $R$, $t$.
  \item Calcule le rapport $W_2/W_1$. Quel montage chauffe le plus ?
\end{enumerate}
\end{exo}

\begin{exo}[chaleur dégagée par deux résistances en série]
Deux résistances $R_1=\qty{10}{\ohm}$ et $R_2=\qty{50}{\ohm}$ sont en série, sous une tension
constante $U=\qty{20}{\volt}$ pendant $t=\qty{9}{\hour}$.
\begin{center}
\begin{circuitikz}[scale=0.85,transform shape]
  \draw (0,0) to[battery1,l_=$U$] (0,2) to[R,l=$R_1$] (2.4,2) to[R,l=$R_2$] (4.8,2) -- (4.8,0) -- (0,0);
\end{circuitikz}
\end{center}
\begin{enumerate}[label=\alph*)]
  \item Calcule la résistance équivalente et l'intensité $I$ (la même partout en série).
  \item En déduire l'énergie thermique totale dégagée (en kJ).
\end{enumerate}
\end{exo}

\begin{exo}[pourquoi transporter à haute tension ?]
Une ligne électrique de résistance $R_{\text{ligne}}=\qty{1}{\ohm}$ transporte une puissance
$P=\qty{1000}{\watt}$.
\begin{enumerate}[label=\alph*)]
  \item Sous une tension de transport $U_1=\qty{100}{\volt}$, calcule l'intensité puis les pertes par
        effet Joule dans la ligne.
  \item On transporte la même puissance sous $U_2=\qty{200}{\volt}$. Recalcule l'intensité et les
        pertes. Par quel facteur ont-elles été divisées ?
  \item Vrai ou faux : « pour réduire les pertes, il faut augmenter l'intensité dans la ligne ».
\end{enumerate}
\end{exo}

\end{document}
